% Options for packages loaded elsewhere
\PassOptionsToPackage{unicode}{hyperref}
\PassOptionsToPackage{hyphens}{url}
\documentclass[
  a4paper,
]{article}
\usepackage{xcolor}
\usepackage[margin=25mm]{geometry}
\usepackage{amsmath,amssymb}
\setcounter{secnumdepth}{-\maxdimen} % remove section numbering
\usepackage{iftex}
\ifPDFTeX
  \usepackage[T1]{fontenc}
  \usepackage[utf8]{inputenc}
  \usepackage{textcomp} % provide euro and other symbols
\else % if luatex or xetex
  \usepackage{unicode-math} % this also loads fontspec
  \defaultfontfeatures{Scale=MatchLowercase}
  \defaultfontfeatures[\rmfamily]{Ligatures=TeX,Scale=1}
\fi
\usepackage{lmodern}
\ifPDFTeX\else
  % xetex/luatex font selection
\fi
% Use upquote if available, for straight quotes in verbatim environments
\IfFileExists{upquote.sty}{\usepackage{upquote}}{}
\IfFileExists{microtype.sty}{% use microtype if available
  \usepackage[]{microtype}
  \UseMicrotypeSet[protrusion]{basicmath} % disable protrusion for tt fonts
}{}
\makeatletter
\@ifundefined{KOMAClassName}{% if non-KOMA class
  \IfFileExists{parskip.sty}{%
    \usepackage{parskip}
  }{% else
    \setlength{\parindent}{0pt}
    \setlength{\parskip}{6pt plus 2pt minus 1pt}}
}{% if KOMA class
  \KOMAoptions{parskip=half}}
\makeatother
\usepackage{longtable,booktabs,array}
\newcounter{none} % for unnumbered tables
\usepackage{calc} % for calculating minipage widths
% Correct order of tables after \paragraph or \subparagraph
\usepackage{etoolbox}
\makeatletter
\patchcmd\longtable{\par}{\if@noskipsec\mbox{}\fi\par}{}{}
\makeatother
% Allow footnotes in longtable head/foot
\IfFileExists{footnotehyper.sty}{\usepackage{footnotehyper}}{\usepackage{footnote}}
\makesavenoteenv{longtable}
\usepackage{graphicx}
\makeatletter
\newsavebox\pandoc@box
\newcommand*\pandocbounded[1]{% scales image to fit in text height/width
  \sbox\pandoc@box{#1}%
  \Gscale@div\@tempa{\textheight}{\dimexpr\ht\pandoc@box+\dp\pandoc@box\relax}%
  \Gscale@div\@tempb{\linewidth}{\wd\pandoc@box}%
  \ifdim\@tempb\p@<\@tempa\p@\let\@tempa\@tempb\fi% select the smaller of both
  \ifdim\@tempa\p@<\p@\scalebox{\@tempa}{\usebox\pandoc@box}%
  \else\usebox{\pandoc@box}%
  \fi%
}
% Set default figure placement to htbp
\def\fps@figure{htbp}
\makeatother
\setlength{\emergencystretch}{3em} % prevent overfull lines
\providecommand{\tightlist}{%
  \setlength{\itemsep}{0pt}\setlength{\parskip}{0pt}}
\usepackage{bookmark}
\IfFileExists{xurl.sty}{\usepackage{xurl}}{} % add URL line breaks if available
\urlstyle{same}
\hypersetup{
  pdftitle={Quadratic Condensation and Relational Local Time in a Many-Wave Complete Relational System --- From Relational Closure without Spacetime as Fundamental Variables to the Kepler Equation and the Inverse-Square Central Force},
  pdfauthor={Noriaki Kihara (WF System Co., Ltd.) ORCID 0009-0004-6753-4020},
  hidelinks,
  pdfcreator={LaTeX via pandoc}}

\title{Quadratic Condensation and Relational Local Time in a Many-Wave
Complete Relational System --- From Relational Closure without Spacetime
as Fundamental Variables to the Kepler Equation and the Inverse-Square
Central Force}
\author{Noriaki Kihara (WF System Co., Ltd.) ORCID 0009-0004-6753-4020}
\date{September 16, 2026 Version DOI 10.5281/zenodo.22788738}

\begin{document}
\maketitle

\subsection{Abstract}\label{abstract}

This study reports an unexpected conservation structure obtained from
numerical experiments that extend the finite complete relational closed
system of our preceding work to a nonlinear system containing many
relational waves. The initial states were generated as linear sums of
two integer harmonic components; the dynamical kernel was never given
the harmonic numbers or the parity classification, and all states were
updated simultaneously by a real antisymmetric generator constructed
only from the relative phases between relational waves sharing a vertex.
Running the preregistered 99 runs for 4096 steps, the anticipated
relaxation-type evolution --- a long stay on the floor, an onset, and
saturation --- did not occur: 5 runs are exact fixed points, and the
remaining 94 depart from the initial plane almost immediately.

On the other hand, across all runs and all times,

\[
\boxed{\sum_m z_m^2=C^2}
\]

was confirmed to be conserved as one and the same complex constant
\(C^2\). This follows rigorously from the fact that the update matrix is
real orthogonal. In this paper we read this identical conservation
structure as \textbf{quadratic condensation}: many internal waves are
represented by a single complex quadratic quantity. For each condensate
\(n\), defining

\[
C_n^2=\sum_m z_{nm}^2,
\]

we have

\[
\sum_n C_n^2=\sum_{n,m}z_{nm}^2,
\]

so that at the higher level one may again impose

\[
\boxed{\sum_n C_n^2=0}.
\]

Quadratic closure therefore closes hierarchically in the same form even
after the internal degrees of freedom are condensed.

Furthermore, the simultaneous conservation of \(H=Z^\dagger Z\) and
\(C^2=Z^TZ\) is equivalent to fixing the Gram matrix of the real and
imaginary parts of the state, which determines the eigenvalues
\(\lambda_\pm=(H\pm|C^2|)/2\) and the shape index
\(\chi=|C^2|/H\in[0,1]\). Under nonvanishing components and
connectivity, \(\chi=1\iff K=0\) holds; in the measurements, exactly the
5 exact fixed points among the 99 runs have \(\chi=1\) (all others have
\(\chi\le0.486\)). For a partition into condensates, an antisymmetric
quadratic flux \(J_{nm}=2Z_n^TK_{nm}Z_m=-J_{mn}\) is defined, with
\(\dot C_n^2=\sum_{m\ne n}J_{nm}\) --- local conservation, pairwise
exchange, and global closure are connected within a single structure. In
addition, for the two-mode initial family, \(\chi\) is determined by a
closed form in \((L,m_a,m_b)\), so the local Gram shape is decided
before any time evolution is executed.

For the minimal two-condensate system

\[
C_1^2+C_2^2=0,
\]

one obtains \(C_2=\pm iC_1\), a 90-degree phase relation. The
transformations preserving this closure can be defined as real
orthogonal transformations, and in the continuous representation as
conservative motion generated by a real antisymmetric generator. No
external physical time needs to be assumed as a fundamental variable.
When two local condensates \(A,B\) exist, B is a clock for A and A is a
clock for B through their relative phase; no absolute clock is required.
A single-valued relational phase \(\tilde U_{AB}\) is determined from
the quadratic representatives alone, and the two lifts on its double
cover give the two time directions \(U_{AB}=\pm i\) at minimal closure.
A complete derivation of time itself is a task for the next stage. What
this paper shows is that a long internal state sequence closes within
one invariant quadratic sector, and that local clocks can be defined as
inter-system relations.

Furthermore, squaring the fixed Gram 2-frame yields an ellipse with one
focus at the origin, whose eccentricity coincides exactly with
\(\chi=|C^2|/H\). In the three-component embedding with the
complementary component \(i\sqrt{|C^2|}\), any nondegenerate elliptic
eccentricity \(0\le e<1\) is realized while exactly preserving quadratic
zero closure, with \(C^2=0\) corresponding to the circular orbit.
Postulating the clock-readout axiom \(dt=r\,d\tau=|u|^2d\tau\), the
Kepler equation for the eccentric anomaly appears, and the same
transformation yields the constancy of the areal velocity and the
inverse-square central force \(\ddot{\mathbf r}=-\mu\,\mathbf r/r^3\).
Thus in this system the central structure of Kepler dynamics is
reconstructed without placing physical time or spatial coordinates as
fundamental variables at the outset.

As for inflation-type relaxation, it is highly likely that the present
99 runs sampled only the two extremes: exact fixed points and
low-symmetry immediate departure. Comparison with the relaxation curves
observed in past high-symmetry 90-degree configurations suggests, as a
candidate, an ``intermediate window'': high symmetry, slightly off the
fixed point, with the centroid near the origin. The verification
procedure is given in the conclusion.

\textbf{Keywords:} complete relational system, quadratic condensation,
quadratic closure, fixed Gram geometry, shape index chi, hierarchical
closure, antisymmetric quadratic flux, relational local time, double
cover, squaring map, Kepler equation, inverse-square central force

\begin{center}\rule{0.5\linewidth}{0.5pt}\end{center}

\section{1. Introduction}\label{introduction}

In the preceding study {[}K1{]}, without placing external time as a
distinguished physical variable from the outset, states at different
phase instants were treated as ordinary state vertices, and the minimal
exactly closed system whose complete network of relations closes was
analyzed. In particular, for

\[
U^K=I
\]

with \(N=1,K=4\), using the base wave and one integer harmonic

\[
S_k=(U^k,U^{mk}),
\]

the orders, complete-relation distances, and quadratic closures of even
versus odd harmonics were compared exactly.

Its direct background, the preceding study {[}K2{]}, placed

\[
M=\frac{N(N-1)}{2}
\]

complex relational waves on the complete graph over \(N\) vertices;
iterating the nonlinear interaction from near a highly symmetric
90-degree phase series, the waves were observed numerically to pass from
the initial cross-shaped band, through a rapid growth of the transverse
component, to a quasi-stable ring of nearly equal amplitudes with widely
dispersed phases.

The original motivation of the present work was simple: to extend the
two-wave structure of paper 1 {[}K1{]} to complete relational systems
with many more relational waves, and to look for general structure among
even/odd harmonics, intrinsic periods, finite cycles, and long-time
evolution.

The outcome differed from the expectation. Even after running the 99
runs to 4096 steps, the intended relaxation-type temporal structure was
not obtained. Instead, a simpler and stronger identity spanning all runs
appeared:

\[
\boxed{\sum_m z_m^2=C^2}.
\]

The purpose of this paper is to organize this result in the following
stages.

\begin{enumerate}
\def\labelenumi{\arabic{enumi}.}
\tightlist
\item
  \textbf{Numerical fact:} In the nonlinear evolution of the 99 runs,
  \(\sum_m z_m^2\) was conserved as one complex constant.
\item
  \textbf{Mathematical consequence:} The many internal waves can be
  exactly contracted to a single quadratic representative \(C^2\), and
  this contraction is closed under hierarchization.
\item
  \textbf{Candidate physical interpretation:} From the relative phase
  relations of multiple condensates, local clocks can possibly be
  defined without assuming an absolute clock.
\item
  \textbf{Geometric consequence:} Kepler ellipses are derived from the
  squaring map of the fixed Gram geometry (eccentricity \(e=\chi\)).
\item
  \textbf{Dynamical consequence:} Under \(dt=|u|^2d\tau\) obtained from
  the clock-readout axiom, the Kepler equation and the inverse-square
  central force are reconstructed.
\end{enumerate}

Items 1, 2, 4, and 5 above are treated as mathematical/numerical
results; item 3 is distinguished as interpretation and as a future
verification task.

\begin{center}\rule{0.5\linewidth}{0.5pt}\end{center}

\section{2. Numerical experimental
system}\label{numerical-experimental-system}

\subsection{2.1 Initial states}\label{initial-states}

For a finite cyclic length \(L\), set

\[
U=\exp\left(\frac{2\pi i}{L}\right).
\]

Let the undirected edges of the complete graph be \(e=(j,k)\), \(j<k\),
and use

\[
\Delta=k-j.
\]

For two mode numbers \(m_a,m_b\), the initial state is

\[
z_{jk}(0)
=
c_a U^{m_a\Delta}
+
c_b U^{m_b\Delta},
\]

where

\[
c_a=e^{-i\pi/(2L)},
\qquad
c_b=e^{+i\pi/(2L)}.
\]

The initial state is thus a \textbf{simple superposition of just two
finite phase waves}. The mode numbers are used only for generating the
initial condition and are never passed to the subsequent dynamics.

\subsection{2.2 Complete relational
dynamics}\label{complete-relational-dynamics}

Two edges \(e,f\) are adjacent when they share a vertex; let \(A_{ef}\)
be the adjacency matrix of the line graph. Writing each state as

\[
z_e=|z_e|e^{i\theta_e},
\]

construct the real antisymmetric generator

\[
K_{ef}(Z)
=
A_{ef}\sin(\theta_f-\theta_e).
\]

The one-step update is

\[
\boxed{
Z'=\exp(\Delta\tau K(Z))Z
}.
\]

Since

\[
K^T=-K,
\]

the matrix

\[
R(Z)=\exp(\Delta\tau K(Z))
\]

is a real orthogonal matrix at every step.

\(K\) is state dependent, so the dynamics as a whole is nonlinear, but
\textbf{each simultaneous update itself is a real orthogonal
transformation}.

\subsection{2.3 The 99-run survey}\label{the-99-run-survey}

The preregistered survey consists of 99 runs in total, all storing the
4097 states

\[
t=0,\ldots,4096.
\]

\textbf{Table 1: Composition of the 99 runs}

{\def\LTcaptype{none} % do not increment counter
\begin{longtable}[]{@{}
  >{\raggedright\arraybackslash}p{(\linewidth - 8\tabcolsep) * \real{0.1875}}
  >{\raggedright\arraybackslash}p{(\linewidth - 8\tabcolsep) * \real{0.1875}}
  >{\raggedright\arraybackslash}p{(\linewidth - 8\tabcolsep) * \real{0.1875}}
  >{\raggedright\arraybackslash}p{(\linewidth - 8\tabcolsep) * \real{0.1875}}
  >{\raggedleft\arraybackslash}p{(\linewidth - 8\tabcolsep) * \real{0.2500}}@{}}
\toprule\noalign{}
\begin{minipage}[b]{\linewidth}\raggedright
Series
\end{minipage} & \begin{minipage}[b]{\linewidth}\raggedright
\(L\)
\end{minipage} & \begin{minipage}[b]{\linewidth}\raggedright
Mode pairs
\end{minipage} & \begin{minipage}[b]{\linewidth}\raggedright
den
\end{minipage} & \begin{minipage}[b]{\linewidth}\raggedleft
Runs
\end{minipage} \\
\midrule\noalign{}
\endhead
\bottomrule\noalign{}
\endlastfoot
Main survey & 8,10,12,16,20,24,32,40 &
(1,3),(1,5),(3,5),(2,4),(2,6),(4,6),(1,2),(2,3),(3,6) & \(L\) & 72 \\
gcd series & 12,16,20,24,32,40 & (3,9),(4,8),(5,10) & \(L\) & 18 \\
Step-size controls & 12,24,32 & (1,3),(2,4),(2,3) & 40 & 9 \\
\textbf{Total} & & & & \textbf{99} \\
\end{longtable}
}

All 99 runs passed QA, and the raw data were never modified by
post-processing.

\begin{center}\rule{0.5\linewidth}{0.5pt}\end{center}

\section{3. Catalog of initial states: the geometry generated by the
two-wave
generator}\label{catalog-of-initial-states-the-geometry-generated-by-the-two-wave-generator}

The catalog (Figure 1) is the completed figure that overlays, for the
same 99 initial states, the point configuration, the two-component
phasor composition paths with multiplicities, the centroid, and the
quadratic representative. The subsections below correspond to the
elements of Figure 1.

\textbf{Figure 1a. Catalog of initial states (completed figure): main
survey, 72 runs}

\begin{figure}
\centering
\pandocbounded{\includegraphics[keepaspectratio,alt={Figure 1a initial states catalog, main 72}]{figs/fig_initial_states_all99_v4_page1_main72_en_20260916.png}}
\caption{Figure 1a initial states catalog, main 72}
\end{figure}

\textbf{Figure 1b. Catalog of initial states (completed figure): gcd
series 18 + den=40 controls 9}

\begin{figure}
\centering
\pandocbounded{\includegraphics[keepaspectratio,alt={Figure 1b initial states catalog, gcd+control}]{figs/fig_initial_states_all99_v4_page2_gcd18_ctrl9_en_20260916.png}}
\caption{Figure 1b initial states catalog, gcd+control}
\end{figure}

\subsection{3.1 Point configuration}\label{point-configuration}

The initial state depends only on \(\Delta=k-j\), so many edges
degenerate to the same point. All edges are nonzero, and as \(L\) grows,
characteristic configurations --- stars, crosses, sparse layouts, double
rings --- appear according to the mode pair.

\subsection{3.2 Two-component composition paths and
multiplicities}\label{two-component-composition-paths-and-multiplicities}

The composition paths in Figure 1 (solid = base-wave component, dotted =
harmonic component) make explicit that each initial complex number
\(z_e(0)\) is generated as the sum of two mode components. However
complicated the dynamics, the starting point is a simple two-component
superposition.

\subsection{3.3 Centroid}\label{centroid}

In all 99 runs the centroid

\[
G=\sum_e z_e
\]

was nonzero. The centroid is not conserved and reorganizes substantially
during the evolution. This paper therefore does not regard the centroid
as a conserved representative of a condensate but treats it as an
\textbf{auxiliary quantity reading the degree of symmetric
configuration}.

\subsection{3.4 Quadratic quantity}\label{quadratic-quantity}

At the time the frame marks of Figure 1 (real and imaginary parts of
\(\sum z^2\)) were created, we called \(\sum z^2\) the
``quadratic-closure residual'' and read it as a deviation from zero. The
subsequent comparison of the full time series showed this reading to be
insufficient.

That

\[
\sum_e z_e^2\neq0
\]

for all 99 runs is not by itself a ``violation of closure''. Rather, for
each run a

\[
\boxed{
\sum_e z_e^2=C^2
}
\]

with a \textbf{run-specific complex constant \(C^2\)} is selected, and
the same value is conserved for all time.

Hereafter we therefore call this the \textbf{quadratic representative}
rather than a residual.

\begin{center}\rule{0.5\linewidth}{0.5pt}\end{center}

\section{4. Complete linear superposition of
waveforms}\label{complete-linear-superposition-of-waveforms}

Linearly superposing, over all \(M\) edges, the two generating waves
carried by each edge shows directly which frequency components the
initial state contains and with what gain.

\textbf{Figure 2a. Complete linear superposition of all \(M\) edge
waves: main 72}

\begin{figure}
\centering
\pandocbounded{\includegraphics[keepaspectratio,alt={Figure 2a linear superposition, main 72}]{figs/fig_waveforms_linear_sum_all99_page1_main72_en_20260916.png}}
\caption{Figure 2a linear superposition, main 72}
\end{figure}

\textbf{Figure 2b. Complete linear superposition of all \(M\) edge
waves: gcd + controls}

\begin{figure}
\centering
\pandocbounded{\includegraphics[keepaspectratio,alt={Figure 2b linear superposition, gcd+control}]{figs/fig_waveforms_linear_sum_all99_page2_gcd18_ctrl9_en_20260916.png}}
\caption{Figure 2b linear superposition, gcd+control}
\end{figure}

With the initial modes \(m_a,m_b\), the total superposition can be
written

\[
W(x)
=
\Re\left[
\left(\sum_e A_e\right)e^{im_ax}
+
\left(\sum_e B_e\right)e^{im_bx}
\right].
\]

For small \(m/L\) the base-wave side often dominates, but in a finite
cyclic system the folding of \(\sin(\pi m/L)\) produces gain inversions
unlike a simple continuous frequency ratio.

The initial configuration is two waves, but the subsequent nonlinear
dynamics knows nothing of the mode labels. Hence the quadratic
conservation obtained later is not the result of ``conserving with the
harmonic numbers held fixed''.

\begin{center}\rule{0.5\linewidth}{0.5pt}\end{center}

\section{5. Final states: internal reorganization after 4096
steps}\label{final-states-internal-reorganization-after-4096-steps}

\textbf{Figure 3a. Final states \(t=4096\): main 72}

\begin{figure}
\centering
\pandocbounded{\includegraphics[keepaspectratio,alt={Figure 3a final states, main 72}]{figs/fig_final_states_all99_page1_main72_en_20260916.png}}
\caption{Figure 3a final states, main 72}
\end{figure}

\textbf{Figure 3b. Final states \(t=4096\): gcd + controls}

\begin{figure}
\centering
\pandocbounded{\includegraphics[keepaspectratio,alt={Figure 3b final states, gcd+control}]{figs/fig_final_states_all99_page2_gcd18_ctrl9_en_20260916.png}}
\caption{Figure 3b final states, gcd+control}
\end{figure}

In most runs with \(L\ge16\), the initially degenerate point
configuration breaks up and reorganizes into ring-like distributions of
nearly equal amplitude. On the other hand, the 5 runs satisfying

\[
m_a+m_b\equiv0\pmod L
\]

have initial states that are real vectors up to a common phase and
remain exact fixed points with

\[
K=0,
\]

preserving the initial configuration. These 5 runs are instances of the
theorem of §6.3 (\(\chi=1\iff K=0\)) and satisfy \(\chi=1\) exactly
(Figure 5).

The internal state may therefore either stay still or reorganize
substantially depending on the run. As the next section shows, in either
case the quadratic representative \(C^2\) is conserved in the same
sense.

\begin{center}\rule{0.5\linewidth}{0.5pt}\end{center}

\section{\texorpdfstring{6. Central result: \(\sum z_m^2=C^2\) closes
exactly for all
time}{6. Central result: \textbackslash sum z\_m\^{}2=C\^{}2 closes exactly for all time}}\label{central-result-sum-z_m2c2-closes-exactly-for-all-time}

\subsection{6.1 Numerical observation}\label{numerical-observation}

\textbf{Figure 4a. Time series of \(\sum z^2\) and centroid \(\sum z\):
main 72}

\begin{figure}
\centering
\pandocbounded{\includegraphics[keepaspectratio,alt={Figure 4a residual time series, main 72}]{figs/fig_residual_timeseries_all99_page1_main72_en_20260916.png}}
\caption{Figure 4a residual time series, main 72}
\end{figure}

\textbf{Figure 4b. Time series of \(\sum z^2\) and centroid \(\sum z\):
gcd + controls}

\begin{figure}
\centering
\pandocbounded{\includegraphics[keepaspectratio,alt={Figure 4b residual time series, gcd+control}]{figs/fig_residual_timeseries_all99_page2_gcd18_ctrl9_en_20260916.png}}
\caption{Figure 4b residual time series, gcd+control}
\end{figure}

In Figure 4, while the centroid \(\sum z\) oscillates strongly and
sometimes amplifies beyond its initial value,

\[
Q_2(t)=\sum_m z_m(t)^2
\]

is a horizontal line in every panel. In the numerical audit over all 99
runs with 4097 states each, the drift of \(Q_2\) stayed within
\(\max_t|Q_2(t)-Q_2(0)|=4.1\times10^{-10}\) (relative \(\sim10^{-12}\)
against \(H_0\sim10^3\)).

What matters at this point is not the numerical approximation

\[
Q_2(t)\approx Q_2(0)
\]

but the following algebraic identity.

\subsection{6.2 Derivation of the exact
conservation}\label{derivation-of-the-exact-conservation}

Write one step as

\[
Z'=R(Z)Z,
\qquad
R(Z)=e^{\Delta\tau K(Z)}.
\]

Since \(K^T=-K\),

\[
R^TR=I.
\]

Therefore

\[
Z'^T Z'
=
Z^T R^T R Z
=
Z^T Z,
\]

that is,

\[
\boxed{
\sum_m z_m'^2
=
\sum_m z_m^2
}.
\]

Since \(R\) is a real matrix, \(R^\dagger=R^T\), and simultaneously

\[
\boxed{
Z'^\dagger Z'
=
Z^\dagger R^TRZ
=
Z^\dagger Z
=
H
}
\]

also holds. Hence

\[
\boxed{H=\text{const},\qquad C^2=\text{const}}
\]

is complete within this section alone.

Defining at the initial value

\[
C^2:=\sum_m z_m(0)^2,
\]

for any step \(t\),

\[
\boxed{
\sum_m z_m(t)^2=C^2
}
\]

holds.

The dynamics is nonlinear through the state dependence of \(K(Z)\), but
the nonlinearity never moves the orbit to another quadratic sector. The
entire orbit closes inside

\[
\mathcal M_{C^2}
=
\left\{
Z\in\mathbb C^M
\mid
Z^T Z=C^2
\right\}.
\]

\subsection{\texorpdfstring{6.3 Fixed Gram geometry and the shape index
\(\chi\)}{6.3 Fixed Gram geometry and the shape index \textbackslash chi}}\label{fixed-gram-geometry-and-the-shape-index-chi}

Decomposing \(Z=x+iy\) (\(x,y\in\mathbb R^M\)),

\[
H=Z^\dagger Z=x^2+y^2,
\qquad
C^2=Z^TZ=x^2-y^2+2i\,x\cdot y,
\]

and the simultaneous conservation of \(H\) and \(C^2\) is equivalent to

\[
x^2=\frac{H+\Re C^2}{2},
\qquad
y^2=\frac{H-\Re C^2}{2},
\qquad
x\cdot y=\frac{\Im C^2}{2},
\]

that is, to fixing the Gram matrix of the real and imaginary parts,

\[
G=
\begin{pmatrix}
x^2 & x\cdot y\\
x\cdot y & y^2
\end{pmatrix}.
\]

The real orthogonal update \(Z'=RZ\) acts as \(x'=Rx,\ y'=Ry\), so this
is a derivation, not an additional assumption. The eigenvalues of \(G\)
are

\[
\boxed{\lambda_\pm=\frac{H\pm|C^2|}{2}}
\]

(since \((\lambda_+-\lambda_-)^2=(x^2-y^2)^2+4(x\cdot y)^2=|C^2|^2\)).
Defining the shape index

\[
\boxed{\chi=\frac{|C^2|}{H}},
\qquad
0\le\chi\le1,
\]

\(\chi=0\) corresponds to a circular 2-frame (equal lengths,
orthogonal), \(0<\chi<1\) to an elliptic 2-frame, and \(\chi=1\) to
degeneration onto a line. The motion therefore does not change the
shape: it \textbf{rotates the fixed-Gram 2-frame inside
\(\mathbb R^M\)}.

The square of the Gram area \(A_{\rm Gram}\) is

\[
\boxed{
A_{\rm Gram}^2=\det G=\lambda_+\lambda_-=\frac{H^2-|C^2|^2}{4}=\frac{H^2}{4}\left(1-\chi^2\right)
},
\]

with \(\chi=1\iff\det G=0\): the frozen class is the state whose
internal 2-frame area has degenerated to zero, while \(\chi=0\) is the
maximal area at fixed \(H\).

\textbf{Theorem (\(\chi\)=1 \(\iff\) K=0).} If all components are
nonzero and the interaction graph (line graph) is connected, then

\[
\boxed{\chi=1\iff K=0}.
\]

Proof: equality in \(|C^2|=|\sum_m z_m^2|\le\sum_m|z_m|^2=H\) holds only
when all \(z_m^2\) share the same phase, i.e.,
\(\theta_m\equiv\theta_0\pmod\pi\), so all relative-phase sines vanish
and \(K=0\). Conversely, if \(K=0\), adjacent phase differences are
\(0\pmod\pi\) and propagate through the whole graph by connectivity.
\(\blacksquare\)

Measurement (Figure 5): among the 99 runs, \(\chi=1\) (to machine
precision) occurs exactly for the 5 frozen-class runs, in agreement with
the analytic fact that for \(m_a+m_b\equiv0\pmod L\), \(z(0)\) is real
times a global phase. The remaining 94 runs lie in
\(\chi\in[0.057,\,0.486]\), with a clear gap below \(\chi=1\). By the
conservation of \(H\) and \(C^2\), \(\chi\) is a \textbf{conserved
label} of each run; its time variation in the invariants series stayed
within \(1.3\times10^{-13}\) over all 4097 steps.

For the two-mode initial family, \(\chi\) is given in closed form in
\((L,m_a,m_b)\). With \(\alpha=\pi/L\), \(M=L(L-1)/2\), and

\[
S_L(r)=\sum_{\Delta=1}^{L-1}(L-\Delta)U^{r\Delta}
=
\begin{cases}
M, & r\equiv0\pmod L,\\[6pt]
\dfrac{LU^r}{1-U^r}
=-\dfrac L2+\dfrac{iL}{2}\cot\dfrac{\pi r}{L},
& r\not\equiv0\pmod L,
\end{cases}
\]

we have

\[
\boxed{
C^2
=
e^{-i\alpha}S_L(2m_a)
+
e^{i\alpha}S_L(2m_b)
+
2S_L(m_a+m_b)
},
\]

\[
\boxed{
H
=
2M+2\Re\!\left[e^{-i\alpha}S_L(m_a-m_b)\right]
},
\qquad
\boxed{
\chi(L,m_a,m_b)=\frac{|C^2|}{H}
}
\]

(derivation in Appendix C). This closed form agrees with the measured
\(\chi,\ C^2,\ H\) of all 99 runs to a maximal relative deviation of
\(9.3\times10^{-15}\), and reproduces the measured maximum
\(\chi(8,3,6)=0.486422965736\ldots\). For the frozen class
\(m_a+m_b\equiv0\pmod L\), the closed form gives \(C^2=H\) (real
positive), so \(\chi=1\) follows analytically. Since \(\chi\) is
conserved, \textbf{the local Gram shape is decided by the initial
integers alone, before any time evolution is executed}.

\textbf{Figure 5. Catalog of the shape index \(\chi\) (all 99 runs;
conserved)}

\begin{figure}
\centering
\pandocbounded{\includegraphics[keepaspectratio,alt={Figure 5 chi catalog}]{figs/fig_chi_all99_en_20260916.png}}
\caption{Figure 5 chi catalog}
\end{figure}

\subsection{6.4 Local state space}\label{local-state-space}

The simultaneous fixing of \(H\) and \(C^2\) defines the internal state
space of one condensate:

\[
\boxed{
\mathcal L_{H,C^2}
=
\{Z\in\mathbb C^M \mid
Z^\dagger Z=H,\;
Z^TZ=C^2\}
}.
\]

For generic rank 2 (\(\chi<1\)) this is the orbit of fixed-Gram
2-frames,

\[
\boxed{
\mathcal L_{H,C^2}
\simeq
O(M)/O(M-2)
},
\qquad
\dim_{\mathbb R}=2M-3.
\]

In the degenerate case (\(\chi=1\), rank 1), the dimension drops to
\(M-1\). Hence

\[
\boxed{
C^2=\text{label of the local closure},
\qquad
\mathcal L_{H,C^2}=\text{its internal local state space}
}.
\]

The orbit \(Z(0)\to Z(1)\to\cdots\) is a motion inside this local space.

\subsection{6.5 Embedding into zero
closure}\label{embedding-into-zero-closure}

For any nonzero \(C^2\),

\[
(iC)^2=-C^2,
\]

so we may write

\[
\boxed{
\sum_m z_m^2+(iC)^2=0
}.
\]

A nonzero quadratic representative therefore does not contradict the
zero-closure principle. Rather, combining the many internal waves
\(\{z_m\}\) with the complementary component \(iC\) that represents them
embeds the completed local condensate exactly into the same quadratic
zero-closure form.

\begin{center}\rule{0.5\linewidth}{0.5pt}\end{center}

\section{7. Quadratic condensation: one complex quantity represents many
waves}\label{quadratic-condensation-one-complex-quantity-represents-many-waves}

\subsection{7.1 Definition}\label{definition}

In this paper, a subsystem in which

\[
C^2=\sum_m z_m^2
\]

holds exactly is called a \textbf{condensate} in the quadratic sense.

This contraction is an identity; it uses no averaging, no low-rank SVD
approximation, no normalization, and no fitting.

With respect to the quadratic quantity, the many internal degrees of
freedom

\[
\{z_1,z_2,\ldots,z_M\}
\]

are \textbf{exactly represented} by the single complex quantity

\[
\boxed{C^2}.
\]

\(C\) itself carries the square-root two-valuedness

\[
C\leftrightarrow -C,
\]

so the fundamental invariant is \(C^2\).

\subsection{7.2 The long internal evolution belongs to the same
condensate}\label{the-long-internal-evolution-belongs-to-the-same-condensate}

Each run of this experiment carries the long state sequence

\[
Z(0),Z(1),\ldots,Z(4096).
\]

The internal configuration changes substantially, but

\[
Z(t)^T Z(t)=C^2
\]

does not. From the higher-level quadratic readout,

\[
\boxed{
Z(0),Z(1),\ldots,Z(4096)
\quad\text{all belong to the same }C^2
}.
\]

This does not mean the 4096 states are identical --- the internal states
differ. But their entire history closes inside one invariant quadratic
sector.

In this sense the experiment shows that

\[
\boxed{
\text{a long local state evolution can close inside one condensate}
}.
\]

\begin{center}\rule{0.5\linewidth}{0.5pt}\end{center}

\section{8. Hierarchical closure and inter-condensate
flux}\label{hierarchical-closure-and-inter-condensate-flux}

\subsection{8.1 Form invariance of hierarchical
contraction}\label{form-invariance-of-hierarchical-contraction}

Consider multiple condensates. Let the internal waves of condensate
\(n\) be \(z_{nm}\) and define

\[
\boxed{
C_n^2=\sum_m z_{nm}^2
}.
\]

By the simple associativity of sums,

\[
\sum_n C_n^2
=
\sum_n\sum_m z_{nm}^2
=
\sum_{n,m}z_{nm}^2.
\]

Imposing at the higher level

\[
\boxed{
\sum_n C_n^2=0
}
\]

therefore automatically gives

\[
\boxed{
\sum_{n,m}z_{nm}^2=0
}.
\]

Grouping several \(C_n\) further,

\[
D_a^2=\sum_{n\in a}C_n^2
\quad\Longrightarrow\quad
\sum_aD_a^2=\sum_nC_n^2.
\]

Raising the hierarchy

\[
z_{nm}\rightarrow C_n\rightarrow D_a\rightarrow\cdots
\]

never changes the form of the quadratic closure.

The central algebraic result of this paper is summarized as

\[
\boxed{
\text{quadratic closure is form-invariant under coarse-graining by quadratic condensation}
}.
\]

\subsection{8.2 Antisymmetric quadratic flux between
condensates}\label{antisymmetric-quadratic-flux-between-condensates}

Partition the full state by condensate, \(Z=(Z_1,Z_2,\ldots)\), and
block the real antisymmetric generator (\(K_{mn}=-K_{nm}^T\)). In the
continuous representation \(\dot Z=KZ\) (the flow of the per-step frozen
\(K\), the same convention as §9.2), with \(q_n=Z_n^TZ_n=C_n^2\),

\[
\dot q_n
=
2Z_n^T\dot Z_n
=
2\sum_m Z_n^TK_{nm}Z_m,
\]

and the self term vanishes by antisymmetry, \(Z_n^TK_{nn}Z_n=0\).
Defining

\[
\boxed{J_{nm}=2Z_n^TK_{nm}Z_m},
\qquad
\boxed{J_{mn}=-J_{nm}},
\]

we obtain

\[
\boxed{
\dot C_n^2
=
\sum_{m\ne n}J_{nm}
},
\qquad
\frac{d}{ds}\sum_nC_n^2=0.
\]

Internal motion never changes its own \(C_n^2\); only relations to other
condensates change \(C_n^2\), and the exchange is antisymmetric. That
is,

\[
\boxed{
\text{local conservation}
+
\text{pairwise exchange}
=
\text{global closure}
}
\]

holds as it stands.

This exchange law is not restricted to the continuous representation.
Freezing \(K_t=K(Z_t)\) at step \(t\) and setting \(Z(s)=e^{sK_t}Z_t\)
(\(0\le s\le\Delta\tau\)), for

\[
\mathcal J_{nm}^{(t)}
=
\int_0^{\Delta\tau}
2Z_n(s)^T K_{nm}^{(t)}Z_m(s)\,ds
\]

we have exactly

\[
\boxed{
C_n^2(t+1)-C_n^2(t)
=
\sum_{m\ne n}\mathcal J_{nm}^{(t)}
},
\qquad
\boxed{
\mathcal J_{mn}^{(t)}=-\mathcal J_{nm}^{(t)}
}.
\]

The antisymmetric quadratic flux is therefore not an approximation of
the continuous representation but an exact exchange law of the
simultaneous update algorithm itself.

\begin{center}\rule{0.5\linewidth}{0.5pt}\end{center}

\section{9. Minimal two-condensate system and conservative motion
without assumed
time}\label{minimal-two-condensate-system-and-conservative-motion-without-assumed-time}

\subsection{9.1 Zero closure of two
condensates}\label{zero-closure-of-two-condensates}

As the minimal higher-level closure, consider

\[
\boxed{
C_1^2+C_2^2=0
}.
\]

The nontrivial solutions are

\[
C_2=\pm iC_1.
\]

The two condensates therefore have

\begin{itemize}
\tightlist
\item
  equal amplitude,
\item
  a 90-degree phase difference.
\end{itemize}

This is the minimal relation of the same type as the 90-degree phase
skeleton that was numerically important in the preceding study {[}K2{]}.
Whether this agreement stems from the same mechanism is to be tested by
the explicit construction of the two-condensate system indicated in the
conclusion.

\subsection{9.2 Transformations preserving the
closure}\label{transformations-preserving-the-closure}

For the state vector

\[
\mathbf C=
\begin{pmatrix}
C_1\\
C_2
\end{pmatrix},
\]

set

\[
\mathbf C'=R\mathbf C.
\]

If

\[
R^TR=I,
\]

then

\[
{\mathbf C'}^T\mathbf C'
=
\mathbf C^TR^TR\mathbf C
=
\mathbf C^T\mathbf C
=
0.
\]

A class of state transformations that do not break the quadratic closure
is thus defined.

Introducing a continuous parameter \(s\) merely as an orbit parameter,

\[
R(s)=e^{sK},
\qquad K^T=-K,
\]

and

\[
\frac{d\mathbf C}{ds}=K\mathbf C,
\]

with

\[
\frac{d}{ds}
\left(
\mathbf C^T\mathbf C
\right)
=
0.
\]

What matters is that \(s\) need not be assumed to be physical absolute
time from the outset. First there is

\[
\boxed{
\text{closure-preserving motion between states}
},
\]

and \(s\) is used only to describe its orbit.

\begin{center}\rule{0.5\linewidth}{0.5pt}\end{center}

\section{10. Local clocks: pairwise relations, not an absolute
clock}\label{local-clocks-pairwise-relations-not-an-absolute-clock}

Even if each condensate carries internal state evolution, no absolute
clock needs to be attached to any single condensate.

Write the quadratic representative of each condensate as

\[
C_n^2=\rho_n e^{i\Theta_n},
\qquad
\rho_n=|C_n^2|,
\qquad
\Theta_n\in\mathbb R/2\pi\mathbb Z.
\]

\(\Theta_n=\arg C_n^2\) is single-valued (the phase of \(C_n\) itself
carries the two-valuedness \(C\leftrightarrow-C\), §7.1).

Define the relational phase from A to B, single-valued from the
quadratic representatives alone:

\[
\boxed{
\tilde U_{AB}
=
e^{i(\Theta_B-\Theta_A)}
=
\frac{C_B^2/|C_B^2|}{C_A^2/|C_A^2|}
}.
\]

The reverse relation is

\[
\boxed{
\tilde U_{BA}=\tilde U_{AB}^{-1}
}.
\]

Under a common phase shift of all condensates,
\(\Theta_n\to\Theta_n+\beta\), \(\tilde U_{AB}\) is invariant: the
absolute phase disappears completely from the relational quantity.

For three systems, the composition law

\[
\boxed{
\tilde U_{AB}\tilde U_{BC}=\tilde U_{AC}
},
\qquad
\boxed{
\tilde U_{AB}\tilde U_{BC}\tilde U_{CA}=1
}
\]

holds immediately. If A reads B and B reads C, the clock relation from A
to C composes. Local clocks are not merely pairwise ratios; they compose
consistently over the complete relational network.

The directed relation of local clocks lives on the double cover of this
single-valued relational quantity. Taking a lift \(U_{AB}\) with

\[
\boxed{
U_{AB}^2=\tilde U_{AB}
},
\]

the reverse is

\[
\boxed{
U_{BA}=U_{AB}^{-1}
}.
\]

At minimal closure \(C_B^2=-C_A^2\),

\[
\tilde U_{AB}=-1,
\]

whose two lifts are \(U_{AB}=+i,\,-i\); reversing the order of the
relation exchanges \(+i\leftrightarrow-i\). In the minimal
two-condensate system, therefore, \textbf{the \(\mathbb Z_2\) fiber of
the double cover itself carries the direction of the local clock}. This
is not a direction variable added from outside; it is intrinsic to the
square-root structure that quadratic condensation possesses from the
start.

No third-party absolute clock is needed here:

\[
\boxed{
\text{B is a clock for A, and A is a clock for B}
}.
\]

In this reading, time appears not as an absolute quantity attached to
one system but as a directed relation between two locally closed
systems.

Viewing the whole as one closed relational structure, there is only

\[
\sum_n C_n^2=0,
\]

and no single absolute clock or absolute time origin needs to be added.

Between local systems, on the other hand, the directed relations

\[
A-B,\qquad B-A
\]

exist naturally.

What follows directly from this structure is:

\[
\boxed{
\text{local clocks can be defined as inter-system relations, without assuming an absolute clock}
}.
\]

\begin{center}\rule{0.5\linewidth}{0.5pt}\end{center}

\section{11. Reconstruction of Kepler dynamics from relational
closure}\label{reconstruction-of-kepler-dynamics-from-relational-closure}

This section shows that combining the fixed Gram geometry of §6.3 and
the zero-closure embedding of §6.5 with one clock-readout axiom
\(dt=r\,d\tau\) (§11.3) reconstructs the Kepler ellipse, the Kepler
equation, the area law, and the inverse-square central force from a
system that does not place physical time or spatial coordinates as
fundamental variables.

\subsection{11.1 Principal-axis representation of the fixed Gram
2-frame}\label{principal-axis-representation-of-the-fixed-gram-2-frame}

For the eigenvalues \(\lambda_\pm=(H\pm|C^2|)/2\) of §6.3, set

\[
a^2=\frac{H+|C^2|}{2},
\qquad
b^2=\frac{H-|C^2|}{2}.
\]

Aligning the global phase to \(\arg C^2=0\) (a gauge choice that changes
neither \(H\) nor \(|C^2|\)), the 2-frame takes principal-axis form.
Using the complementary component of §6.5 directly, with
\(c_0:=\sqrt{|C^2|}\), take the three-component state

\[
\boxed{
\mathbf V(0)=
\begin{pmatrix}
a\\ ib\\ ic_0
\end{pmatrix}
}.
\]

Then

\[
V_1^2+V_2^2+V_3^2=a^2-b^2-c_0^2=|C^2|-|C^2|=0,
\]

so \textbf{exact quadratic zero closure} holds. Taking the real
antisymmetric generator

\[
\boxed{
K_3=
\begin{pmatrix}
0&\omega&0\\
-\omega&0&0\\
0&0&0
\end{pmatrix}
},
\qquad K_3^T=-K_3,
\]

the complementary component \(V_3=ic_0\) stays fixed while the first
component of the closure-preserving motion
\(\mathbf V(\tau)=e^{\tau K_3}\mathbf V(0)\) reads out as

\[
\boxed{
u(\tau)=V_1(\tau)=a\cos(\omega\tau)+ib\sin(\omega\tau)
}.
\]

Here \(\tau\) is not external physical time but the parameter of the
closure-preserving orbit. Directly from \(K_3^T=-K_3\),

\[
u''=-\omega^2u,
\]

and throughout, \(|V_1|^2+|V_2|^2=a^2+b^2=H\),
\(V_1^2+V_2^2=a^2-b^2=|C^2|\), and the total zero closure
\(\sum_iV_i^2=0\) are identically preserved.

In particular, if \(C^2=0\) then \(a=b\), the complementary component
vanishes, and \(e=\chi=0\): \textbf{the minimal two-condensate zero
closure of §9 (\(C_2=\pm iC_1\)) corresponds to the circular Kepler
orbit}. The general elliptic orbit corresponds to \(C^2\neq0\), i.e., to
zero closure with the complementary component \(i\sqrt{|C^2|}\), and any
nondegenerate elliptic eccentricity \(0\le e=\chi<1\) is realized while
exactly preserving quadratic zero closure. \(\chi=1\) is the boundary
where \(b=0\), \(B=0\), and \(w\) degenerates to a segment through the
focus, in agreement with the rank-1 degeneracy of §6.3.

\subsection{11.2 The focal ellipse by the squaring
map}\label{the-focal-ellipse-by-the-squaring-map}

Take the squaring map

\[
\boxed{w=u^2}.
\]

Setting \(E=2\omega\tau\),

\[
w=D+A\cos E+iB\sin E,
\qquad
A=\frac H2,\quad
D=\frac{|C^2|}2,\quad
B=\frac12\sqrt{H^2-|C^2|^2},
\]

and

\[
\boxed{A^2-B^2=D^2}
\]

holds (\(A^2-B^2=|C^2|^2/4=D^2\)). The ellipse therefore has center
\(D\), semi-major axis \(A\), and focal distance
\(c_f=\sqrt{A^2-B^2}=D\), so \textbf{one focus coincides with the
origin}. The eccentricity is

\[
\boxed{
e=\frac DA=\frac{|C^2|}{H}=\chi
}.
\]

\textbf{The shape index \(\chi\) is precisely the eccentricity of the
Kepler ellipse after the squaring map.}

The distance from the focus (origin) is

\[
\boxed{r=|w|=|u|^2},
\qquad
r=A(1+e\cos E).
\]

To match the usual Kepler convention (pericenter origin), set
\(E_K=E-\pi\):

\[
r=A(1-e\cos E_K).
\]

\subsection{11.3 Clock-readout axiom}\label{clock-readout-axiom}

In this paper the readout of the relational clock is posited as one
\textbf{axiom}:

\[
\boxed{dt=r\,d\tau=|u|^2\,d\tau}
\]

(of the same form as a Sundman-type transformation). This paper does not
derive this readout rule.

Note that the hyperbolic local closure \(r^2-c^2t^2=R^2\) yields
\(c\,dt=r\,d\eta\) via \(r=R\cosh\eta\), \(ct=R\sinh\eta\), but the
hyperbolic angle \(\eta\) cannot be identified with the orbit parameter
\(\tau\): on a Kepler orbit \(r(\tau)=A(1+e\cos2\omega\tau)\) is bounded
and periodic, whereas \(\eta=\gamma\tau\) would give the unbounded
\(r=R\cosh(\gamma\tau)\), incompatible with §11.2 (recorded here as part
of the investigation). The hyperbolic closure is therefore kept as a
separate relational closure, and \(\eta\) and \(\tau\) are not
identified. Deriving the readout rule \(dt/d\tau=|u|^2\) from the
relational clock \(\tilde U_{AB}\) or from the quadratic closure
structure is an open problem (§15.2).

\subsection{11.4 The Kepler equation}\label{the-kepler-equation}

From \(E=2\omega\tau\), \(d\tau=dE/(2\omega)\), so

\[
dt=\frac{A}{2\omega}(1+e\cos E)\,dE,
\]

and integrating,

\[
t-t_0=\frac{A}{2\omega}(E+e\sin E).
\]

Moving to the pericenter convention \(E_K=E-\pi\) with pericenter
passage time \(t_p=t_0+\pi A/(2\omega)\),

\[
\boxed{M=E_K-e\sin E_K},
\qquad
M=n(t-t_p),
\qquad
\boxed{n=\frac{2\omega}A}.
\]

In a system that did not place external physical time as a fundamental
variable, the Kepler equation appears under the clock-readout axiom
\(dt=r\,d\tau\).

\subsection{11.5 The area law (Kepler's second
law)}\label{the-area-law-keplers-second-law}

From \(w=u^2\) and \(dt=|u|^2d\tau\),

\[
\dot w=\frac{dw/d\tau}{dt/d\tau}=\frac{2uu'}{|u|^2}=\frac{2u'}{\bar u}.
\]

The specific angular momentum about the focus at the origin,
\(h=\Im(\bar w\dot w)\), satisfies

\[
h=2\Im(\bar uu'),
\qquad
\Im(\bar uu')=ab\omega,
\]

so

\[
\boxed{h=2ab\omega=\mathrm{const}},
\qquad
\boxed{\frac{dA_{\rm swept}}{dt}=\frac h2=ab\omega=\mathrm{const}}.
\]

This is Kepler's second law.

\subsection{11.6 The inverse-square central force and Kepler's third
law}\label{the-inverse-square-central-force-and-keplers-third-law}

Use the oscillator-side conserved quantity
\(E_{\rm osc}=\frac12(|u'|^2+\omega^2|u|^2)\). Differentiating
\(\dot w=2u'/\bar u\) once more with respect to \(t\), and using
\(u''=-\omega^2u\) and \(r=|w|=|u|^2\),

\[
\ddot w
=\frac{2}{|u|^2}\left[\frac{u''}{\bar u}-\frac{u'\bar u'}{\bar u^2}\right]
=-\frac{2}{u\bar u^3}\left(\omega^2u\bar u+u'\bar u'\right)
=-\frac{4E_{\rm osc}}{u\bar u^3},
\]

that is,

\[
\boxed{\ddot w=-\frac{4E_{\rm osc}}{|w|^3}w},
\qquad
\boxed{\ddot{\mathbf r}=-\mu\frac{\mathbf r}{r^3}},
\qquad
\boxed{\mu=4E_{\rm osc}}.
\]

On the elliptic solution,
\(E_{\rm osc}=\frac12\omega^2(a^2+b^2)=\frac12\omega^2H\), so

\[
\boxed{\mu=2\omega^2H}.
\]

Moreover, from \(A=H/2\) and \(n=2\omega/A\),

\[
\boxed{n^2A^3=4\omega^2A=2\omega^2H=\mu},
\]

consistent with Kepler's third law.

All formulas of this section were verified independently both
symbolically and numerically (the focal ellipse form, \(e=\chi\), the
Kepler equation, the constancy of \(h\), the residual of
\(\ddot w+\mu w/r^3\), and \(n^2A^3=\mu\), all within
numerical-differentiation accuracy).

\begin{center}\rule{0.5\linewidth}{0.5pt}\end{center}

\section{12. Inflation-type transitions: the two extremes sampled by
this
survey}\label{inflation-type-transitions-the-two-extremes-sampled-by-this-survey}

\subsection{12.1 The indicator of departure from the initial
plane}\label{the-indicator-of-departure-from-the-initial-plane}

With the same definition as the preceding inflation experiments, let
\(E_0\) be the real two-dimensional plane spanned by the real and
imaginary parts of the initial state \(Z_0\), and let \(Z_\perp(t)\) be
the orthogonal component:

\[
f(t)
=
\frac{\langle Z_\perp,Z_\perp\rangle}
{\langle Z,Z\rangle}.
\]

\textbf{Figure 6a. Inflation indicator \(H_\perp/H\): main 72}

\begin{figure}
\centering
\pandocbounded{\includegraphics[keepaspectratio,alt={Figure 6a Hperp/H, main 72}]{figs/fig_inflation_hperp_all99_page1_main72_en_20260916.png}}
\caption{Figure 6a Hperp/H, main 72}
\end{figure}

\textbf{Figure 6b. Inflation indicator \(H_\perp/H\): gcd + den=40
controls}

\begin{figure}
\centering
\pandocbounded{\includegraphics[keepaspectratio,alt={Figure 6b Hperp/H, gcd+control}]{figs/fig_inflation_hperp_all99_page2_gcd18_ctrl9_en_20260916.png}}
\caption{Figure 6b Hperp/H, gcd+control}
\end{figure}

In all 99 runs, \(f(0)\) lies essentially at the display floor by
construction. Afterwards,

\begin{itemize}
\tightlist
\item
  5 runs remained at their fixed points,
\item
  94 runs reached \(O(1)\) within the first few steps.
\end{itemize}

That is, the relaxation curve seen in the preceding study {[}K2{]},

\[
\text{floor}
\rightarrow
\text{onset}
\rightarrow
\text{saturation},
\]

was not contained in a single one of the present 99 runs.

Note that, from the fixed Gram geometry of §6.3, \(H_\perp/H\) measures
not a growth of energy but \textbf{the degree to which the fixed-Gram
2-frame rotates and develops out of the initial 2-plane}.

\subsection{12.2 Present hypotheses}\label{present-hypotheses}

What this survey establishes is that the 99 runs contain only the two
sides of the relaxation curve (exact fixed points and immediate
saturation). The generation condition is to be narrowed by comparing the
following three regions.

\begin{enumerate}
\def\labelenumi{\arabic{enumi}.}
\tightlist
\item
  \textbf{The over-symmetric extreme:} exact fixed points satisfying the
  conjugate-pair condition.
\item
  \textbf{The low-symmetry extreme:} the 94 runs departing immediately
  from the initial plane.
\item
  \textbf{In between:} the region with a finite floor stay and onset
  observed in past high-symmetry 90-degree systems.
\end{enumerate}

Also, the present fixed points have comparatively small centroids, but
small centroids alone do not freeze a run. What decides freezing is
\(\chi=1\) (the theorem of §6.3); in the measurements too, only the 5
fixed runs have \(\chi=1\) while the other 94 have \(\chi\le0.486\)
(Figure 5). Hence the candidate combination is

\[
\boxed{
\text{high symmetry}
+
\text{centroid near the origin}
+
\text{yet not an exact fixed point}
},
\]

to be verified by the perturbation-series experiments indicated in the
conclusion.

\begin{center}\rule{0.5\linewidth}{0.5pt}\end{center}

\section{13. On the ``arrow of time'' and
entropy}\label{on-the-arrow-of-time-and-entropy}

In the basic structure of this paper, entropy need not be introduced to
define the direction of time.

Given a two-body system \(A,B\), the directed relation

\[
\tilde U_{BA}
=
\tilde U_{AB}^{-1}
\]

exists from the start. The direction of the local clock can therefore be
defined by the pairwise relation.

Describing the whole from outside, on the other hand, there is only the
closed structure

\[
\sum_nC_n^2=0,
\]

and no absolute arrow of time needs to be placed as a fundamental
variable.

If entropy growth is observed under macroscopic coarse-graining, it may
correspond to state reorganization along the local relational clocks.
This paper does not assume

\[
\boxed{
\text{entropy growth creates the direction of time}
};
\]

rather it leaves as a future topic

\[
\boxed{
\text{local relational evolution comes first; the entropy direction may be its coarse-grained image}
}.
\]

This is not a mathematical main result of this paper but a candidate
physical interpretation arising from the structure of quadratic
condensation and local clocks.

\begin{center}\rule{0.5\linewidth}{0.5pt}\end{center}

\section{14. Relation to prior work}\label{relation-to-prior-work}

The quadratic condensation and hierarchical closure of this study were
not derived by assuming the following external works. We first found

\[
\sum_mz_m^2=C^2
\]

from the numerical experiments and the algebraic structure, and only
afterwards examined the conceptual relation to existing studies that do
not take external time as a fundamental variable.

Page and Wootters {[}E1{]} treat the closed total system as a stationary
state and describe observed evolution through correlations with an
internal clock. The common point with our work is that no external
absolute clock is fundamental. In contrast, we do not start from a
heterogeneous clock/system split: homogeneous local condensates A and B
read each other as clocks.

Rovelli {[}E2{]} discussed describing observable evolution even where a
fundamental time variable is not defined. Our study does not start from
the canonical formalism of quantum gravity but from the quadratic
closure of finite complex relational waves and real orthogonal
conservative motion.

Boette and Rossignoli {[}E3{]} treat finite-dimensional quantum clocks
and discrete history states, including multiple instants in one static
structure. The contact point is that a long state sequence can be
treated as the interior of one closed structure; however, we assume no
independent clock Hilbert space.

Wharton {[}E4{]} argued against taking the sequential generation of the
future from initial conditions as the only fundamental form, using
time-symmetric boundary conditions constraining whole histories. The
contact point is that no absolute future-generating direction needs to
be placed on the side of the total system; the difference is that our
starting point is finite quadratic closure and relational phase motion,
not an action principle or continuous-field boundary conditions.

The correspondence between the harmonic oscillator and the Kepler
problem via the squaring map is itself the known correspondence going
back to Bohlin {[}E5{]} and Levi-Civita {[}E6{]}. This study did not
construct the model by assuming that known transformation: the squaring
map appeared independently from relational closure and fixed Gram
geometry, and the comparison with the known correspondence was made
after the results were obtained. The difference is that in our study the
time variable \(t\) itself is not given externally but appears as the
relational clock via the clock-readout axiom \(dt=r\,d\tau\).

\begin{center}\rule{0.5\linewidth}{0.5pt}\end{center}

\section{15. What this paper establishes / does not
establish}\label{what-this-paper-establishes-does-not-establish}

\subsection{15.1 Established mathematical and numerical
results}\label{established-mathematical-and-numerical-results}

\begin{enumerate}
\def\labelenumi{\arabic{enumi}.}
\tightlist
\item
  In the present dynamics \(K^T=-K\), and each step's update
  \(R=e^{\Delta\tau K}\) is real orthogonal.
\item
  Therefore \[
  Z^TZ=\sum_mz_m^2=C^2
  \] is exactly conserved.
\item
  The 99-run, 4096-step data confirmed this conservation within
  numerical precision.
\item
  The many waves contract exactly to one \(C^2\) with respect to the
  quadratic quantity.
\item
  For multiple condensates, with \[
  C_n^2=\sum_mz_{nm}^2,
  \] the identity \[
  \sum_nC_n^2=\sum_{n,m}z_{nm}^2
  \] holds exactly.
\item
  The nontrivial solutions of the two-condensate zero closure \[
  C_1^2+C_2^2=0
  \] have a 90-degree phase difference.
\item
  The closure-preserving motion is defined by real orthogonal
  transformations, and in the continuous representation by real
  antisymmetric generators.
\item
  The simultaneous conservation of \(H\) and \(C^2\) is equivalent to
  fixing the Gram matrix of the real and imaginary parts, with
  eigenvalues \[
  \lambda_\pm=\frac{H\pm|C^2|}{2}.
  \] The shape index \(\chi=|C^2|/H\in[0,1]\) is a conserved label of
  each run.
\item
  Under nonvanishing components and line-graph connectivity,
  \(\chi=1\iff K=0\). In the measurements, \(\chi=1\) occurs exactly for
  the 5 frozen runs (all others have \(\chi\le0.486\), Figure 5).
\item
  For a condensate partition, \(J_{nm}=2Z_n^TK_{nm}Z_m\) satisfies
  \(J_{mn}=-J_{nm}\) and \[
  \dot C_n^2=\sum_{m\ne n}J_{nm},
  \qquad
  \frac{d}{ds}\sum_nC_n^2=0
  \] (the self term vanishes).
\item
  The local state space is, generically (\(\chi<1\)),
  \(\mathcal L_{H,C^2}\simeq O(M)/O(M-2)\) of real dimension \(2M-3\)
  (and \(M-1\) for \(\chi=1\)).
\item
  The relational phase \(\tilde U_{AB}=e^{i(\Theta_B-\Theta_A)}\)
  (\(\Theta_n=\arg C_n^2\)) is single-valued and common-phase invariant,
  and satisfies \[
  \tilde U_{BA}=\tilde U_{AB}^{-1},
  \qquad
  \tilde U_{AB}\tilde U_{BC}=\tilde U_{AC}.
  \] The double-cover lift \(U_{AB}^2=\tilde U_{AB}\) gives
  \(U_{AB}=\pm i\) at minimal closure, with the \(\mathbb Z_2\) fiber
  carrying the direction.
\item
  The \(\chi\) of the two-mode initial family is given by the closed
  form in \((L,m_a,m_b)\) (via \(S_L(r)\), Appendix C) and agrees with
  the measurements of all 99 runs to machine precision.
\item
  Under the squaring map \(w=u^2\), the fixed Gram ellipse maps to an
  ellipse with a focus at the origin whose eccentricity is
  \(e=\chi=|C^2|/H\).
\item
  Under the clock-readout axiom \(dt=r\,d\tau=|u|^2d\tau\), the Kepler
  equation \(M=E_K-e\sin E_K\) (with \(n=2\omega/A\)) is derived (the
  axiom itself is not derived).
\item
  The same transformation yields the constancy of the areal velocity
  \(h=2ab\omega\) and the inverse-square central force
  \(\ddot{\mathbf r}=-\mu\mathbf r/r^3\) (with \(\mu=2\omega^2H\) and
  \(n^2A^3=\mu\)).
\end{enumerate}

\subsection{15.2 Not yet established}\label{not-yet-established}

\begin{enumerate}
\def\labelenumi{\arabic{enumi}.}
\tightlist
\item
  Which actual particles, fields, or spatial regions the \(C_n\)
  correspond to.
\item
  That the 90-degree phase relation is identical to actual spacetime
  structure.
\item
  That the local clock obtained from relational phase agrees with
  physical proper time.
\item
  The necessary and sufficient conditions for inflation-type relaxation
  curves.
\item
  That high symmetry and near-origin centroid conditions guarantee
  relaxation-type evolution.
\item
  That the entropy direction necessarily follows the direction of local
  relational clocks.
\item
  That this model fully reproduces the time concept of general
  relativity or standard quantum theory.
\item
  The derivation of the clock-readout axiom \(dt=r\,d\tau\) from the
  relational clock \(\tilde U\) or from the quadratic closure structure.
\end{enumerate}

\begin{center}\rule{0.5\linewidth}{0.5pt}\end{center}

\section{16. Conclusion}\label{conclusion}

This study extended the minimal two-wave closed system of paper 1
{[}K1{]} to many-wave complete relational dynamics. The anticipated
relaxation-type temporal structure was not obtained in the 99 runs. From
this outcome, a simpler and stronger invariant structure became clear:

\[
\boxed{
\sum_m z_m^2=C^2
}.
\]

Since it can be written as

\[
\sum_mz_m^2+(iC)^2=0,
\]

the many internal waves are exactly gathered, with respect to the
quadratic quantity, into a single complex condensed quantity \(C^2\).

For multiple condensates,

\[
C_n^2=\sum_mz_{nm}^2
\]

can be defined, and

\[
\boxed{
\sum_nC_n^2=0
}
\]

can again be used, in the same form, as the higher-level closure.
Quadratic closure therefore loses no form under condensing the internal
degrees of freedom and raising the hierarchy.

The minimal two-condensate system

\[
C_1^2+C_2^2=0
\]

carries the 90-degree phase relation, and the closure-preserving motion
can be defined without assuming external absolute time. Given two local
condensates A and B, the phase difference from A to B and that from B to
A exist as sign-reversed relations of each other. Hence, without placing
an absolute clock, the relational local clock

\[
\boxed{
\text{B is a clock for A, and A is a clock for B}
}
\]

can be constructed.

Viewing the whole at once, neither an absolute clock nor an absolute
arrow of time needs to be introduced as a fundamental variable; between
local systems there are directed relations from which local order can be
read.

The most compact conclusion of this paper is the chain

\[
\boxed{
K^T=-K
\Longrightarrow
(H,C^2)\ \text{conserved}
\Longrightarrow
\text{fixed-Gram local space}
\Longrightarrow
\text{quadratic condensation and hierarchical closure}
\Longrightarrow
\text{relational local clocks}
\Longrightarrow
w=u^2
\Longrightarrow
e=\chi
\Longrightarrow
dt=r\,d\tau
\Longrightarrow
\text{Kepler equation}
\Longrightarrow
\frac1{r^2}\ \text{central force}
}.
\]

The central result of this paper is that, from finite relational closure
that does not assume physical time and spatial coordinates as
fundamental variables, conserved geometry, local clocks, the Kepler
equation, and the inverse-square central force are connected through one
quadratic structure.

The next-stage task is to construct multiple condensates preserving
quadratic closure explicitly, and to test independently whether the past
inflation-type relaxation curves are reproduced under high-symmetry
90-degree configurations with near-origin centroids.

\begin{center}\rule{0.5\linewidth}{0.5pt}\end{center}

\section{References}\label{references}

\subsection{Self-citations}\label{self-citations}

\textbf{{[}K1{]}} N. Kihara, ``A Finite Complete-Relation Closed System
with Time Evolution as Vertices --- Exact Minimal Closure at \(N=1,K=4\)
under \(U^K=I\) and Complete-Network Analysis of Even and Odd
Harmonics,'' 2026. Version DOI: 10.5281/zenodo.22763330; Concept DOI:
10.5281/zenodo.22763329.

\textbf{{[}K2{]}} N. Kihara, ``Mechanism of Inflation-like Rapid
Expansion in a Self-Consistent Closed System of Relational Waves,''
2026. Version DOI: 10.5281/zenodo.22176949; Concept DOI:
10.5281/zenodo.22112008.

\subsection{External references: conceptual
neighbors}\label{external-references-conceptual-neighbors}

\textbf{{[}E1{]}} D. N. Page and W. K. Wootters, ``Evolution without
evolution: Dynamics described by stationary observables,''
\emph{Physical Review D} \textbf{27}, 2885--2892 (1983). DOI:
10.1103/PhysRevD.27.2885.

\textbf{{[}E2{]}} C. Rovelli, ``Time in quantum gravity: An
hypothesis,'' \emph{Physical Review D} \textbf{43}, 442--456 (1991).
DOI: 10.1103/PhysRevD.43.442.

\textbf{{[}E3{]}} A. Boette and R. Rossignoli, ``History states of
systems and operators,'' \emph{Physical Review A} \textbf{98}, 032108
(2018). DOI: 10.1103/PhysRevA.98.032108.

\textbf{{[}E4{]}} K. Wharton, ``Time-Symmetric Boundary Conditions and
Quantum Foundations,'' \emph{Symmetry} \textbf{2}, 272--283 (2010). DOI:
10.3390/sym2010272.

\textbf{{[}E5{]}} K. Bohlin, ``Note sur le problème des deux corps et
sur une intégration nouvelle dans le problème des trois corps,''
\emph{Bulletin Astronomique} \textbf{28}, 113--119 (1911).

\textbf{{[}E6{]}} T. Levi-Civita, ``Sur la régularisation du problème
des trois corps,'' \emph{Acta Mathematica} \textbf{42}, 99--144 (1920).

\begin{center}\rule{0.5\linewidth}{0.5pt}\end{center}

\section{Appendix A. List of all figure files
used}\label{appendix-a.-list-of-all-figure-files-used}

This paper cites \textbf{only the completed figures} from the figure
catalogs created on 2026-09-16 (the intermediate versions v1--v3 of the
initial-state catalog are not cited). The English-labeled figures below
were generated by faithful copies of the original programs, changing
only label strings and output filenames.

\begin{enumerate}
\def\labelenumi{\arabic{enumi}.}
\tightlist
\item
  \texttt{fig\_initial\_states\_all99\_v4\_page1\_main72\_en\_20260916.png}
\item
  \texttt{fig\_initial\_states\_all99\_v4\_page2\_gcd18\_ctrl9\_en\_20260916.png}
\item
  \texttt{fig\_waveforms\_linear\_sum\_all99\_page1\_main72\_en\_20260916.png}
\item
  \texttt{fig\_waveforms\_linear\_sum\_all99\_page2\_gcd18\_ctrl9\_en\_20260916.png}
\item
  \texttt{fig\_final\_states\_all99\_page1\_main72\_en\_20260916.png}
\item
  \texttt{fig\_final\_states\_all99\_page2\_gcd18\_ctrl9\_en\_20260916.png}
\item
  \texttt{fig\_residual\_timeseries\_all99\_page1\_main72\_en\_20260916.png}
\item
  \texttt{fig\_residual\_timeseries\_all99\_page2\_gcd18\_ctrl9\_en\_20260916.png}
\item
  \texttt{fig\_inflation\_hperp\_all99\_page1\_main72\_en\_20260916.png}
\item
  \texttt{fig\_inflation\_hperp\_all99\_page2\_gcd18\_ctrl9\_en\_20260916.png}
\item
  \texttt{fig\_chi\_all99\_en\_20260916.png}
\end{enumerate}

Figures 1--4 (a/b) correspond to items 1--8, Figures 6a/b to items
9--10, and Figure 5 (the \(\chi\) catalog) to item 11.

\begin{center}\rule{0.5\linewidth}{0.5pt}\end{center}

\section{Appendix B. Reproducibility
notes}\label{appendix-b.-reproducibility-notes}

The 99-run raw-trajectory database of this paper stores 4097 states per
run together with QA, metadata, invariants, and SHA256 records. All 99
runs were fixed after the collection was completed, and the figures of
this paper were generated by read-only access to the stored
\texttt{states.npz}. No trajectory was rerun and no condition was
changed at figure-generation time.

The central conservation law used in this paper,

\[
\sum_mz_m^2=C^2,
\]

is not a hypothesis introduced by reinterpretation after the figures
were made; it is the exact algebraic conservation law inherent in the
real antisymmetric generator and the real orthogonal update. What is new
is the rereading of this conserved quantity not as a ``residual from
zero'' but as the ``condensed quantity representing the many waves''.

\begin{center}\rule{0.5\linewidth}{0.5pt}\end{center}

\section{\texorpdfstring{Appendix C. Closed-form derivation of
\(\chi\)(L, m\_a,
m\_b)}{Appendix C. Closed-form derivation of \textbackslash chi(L, m\_a, m\_b)}}\label{appendix-c.-closed-form-derivation-of-chil-m_a-m_b}

For the initial condition
\(z_{jk}(0)=c_aU^{m_a\Delta}+c_bU^{m_b\Delta}\) (\(c_a=e^{-i\pi/(2L)}\),
\(c_b=e^{+i\pi/(2L)}\), \(\Delta=k-j\)), the number of edges with
difference \(\Delta=1,\ldots,L-1\) is \(L-\Delta\). Define

\[
S_L(r):=\sum_{\Delta=1}^{L-1}(L-\Delta)U^{r\Delta},
\qquad U=e^{2\pi i/L}.
\]

\textbf{Evaluation of \(S_L(r)\).} If \(r\equiv0\pmod L\), then
\(U^{r\Delta}=1\) and \(S_L(r)=\sum_{\Delta=1}^{L-1}(L-\Delta)=M\). If
\(r\not\equiv0\), then for \(\omega=U^r\neq1\) with \(\omega^L=1\),

\[
\sum_{\Delta=1}^{L-1}\omega^\Delta=-1,
\qquad
\sum_{\Delta=1}^{L-1}\Delta\,\omega^\Delta=\frac{L}{\omega-1}
\]

(the latter is the standard sum using \(\omega^L=1\)). Therefore

\[
S_L(r)=L\cdot(-1)-\frac{L}{\omega-1}
=-L\,\frac{\omega}{\omega-1}
=\frac{L\omega}{1-\omega}.
\]

With \(\omega=e^{i\theta}\) (\(\theta=2\pi r/L\)),
\(\dfrac{\omega}{1-\omega}=-\dfrac12+\dfrac i2\cot\dfrac\theta2\), so

\[
S_L(r)=-\frac L2+\frac{iL}2\cot\frac{\pi r}{L}.
\]

\textbf{\(C^2\) and \(H\).} Using \(c_a^2=e^{-i\alpha}\),
\(c_b^2=e^{i\alpha}\), \(c_ac_b=1\), \(c_a\bar c_b=e^{-i\alpha}\)
(\(\alpha=\pi/L\)), and collecting edge sums by \(\Delta\) class,

\[
C^2=\sum_e z_e^2
=e^{-i\alpha}S_L(2m_a)+e^{i\alpha}S_L(2m_b)+2S_L(m_a+m_b),
\]

\[
H=\sum_e|z_e|^2
=2M+e^{-i\alpha}S_L(m_a-m_b)+e^{i\alpha}S_L(m_b-m_a)
=2M+2\Re\!\left[e^{-i\alpha}S_L(m_a-m_b)\right]
\]

(using \(S_L(-r)=\overline{S_L(r)}\)). Hence

\[
\chi(L,m_a,m_b)=\frac{|C^2|}{H}.
\]

\textbf{Frozen class.} If \(m_a+m_b\equiv0\pmod L\), then
\(S_L(m_a+m_b)=M\) and \(m_a-m_b\equiv2m_a\pmod L\), so

\[
C^2=2M+2\Re\!\left[e^{-i\alpha}S_L(2m_a)\right]=H,
\]

i.e., \(C^2\) is real positive and \(\chi=1\) follows analytically.

\textbf{Numerical check.} This closed form agrees with the measured
\(\chi,\ C^2,\ H\) of all 99 runs to a maximal relative deviation of
\(9.3\times10^{-15}\).

\end{document}
